%%%% CAPÍTULO 2 - FUNDAMENTAÇÃO TEÓRICA

\chapter{Fundamentação Teórica}\label{cap:fundamentacao}

Este capítulo apresenta a base teórica necessária para a compreensão dos
conceitos discutidos ao longo do trabalho. Inicialmente, aborda-se o fenômeno do
Fator Casa no contexto esportivo e suas variáveis determinantes. Em seguida,
detalha-se a metodologia CRISP-DM, utilizada como guia estrutural da pesquisa.
Por fim, apresentam-se os modelos estatísticos aplicados a resultados de futebol
e os critérios de avaliação de previsões probabilísticas.

\section{O Fenômeno do Fator Casa}\label{sec:fator-casa}

O conceito de Fator Casa (\textit{home advantage}) é amplamente documentado na
literatura da ciência do esporte e refere-se à vantagem estatística consistente
que os times mandantes possuem sobre os visitantes. Segundo estudos clássicos
estruturados por \citeonline{Courneya1992} e \citeonline{Pollard2006}, essa
vantagem pode ser quantificada pela porcentagem de vitórias ou pontos
conquistados em casa, que historicamente supera a marca de 50\% em diversas
modalidades de esportes coletivos.

\citeonline{Nevill1999} realizaram uma ampla revisão da literatura e apontaram
que o Fator Casa é um fenômeno multifatorial. As principais variáveis que
sustentam esse comportamento incluem:

\begin{itemize}
  \item \textbf{Apoio da torcida:} a presença do público local atua como um
    forte motivador psicológico para os jogadores mandantes, além de criar um
    ambiente de pressão sobre os adversários;
  \item \textbf{Familiaridade com o ambiente:} o conhecimento detalhado das
    dimensões do campo, das condições do gramado e até mesmo do clima local
    confere uma vantagem tática sutil, mas real;
  \item \textbf{Logística e cansaço de viagem:} equipes visitantes enfrentam o
    desgaste físico e mental decorrente de deslocamentos, alterações de rotina e
    hospedagens, fatores que podem impactar o rendimento atlético;
  \item \textbf{Viés de arbitragem:} estudos apontam que a pressão exercida pela
    torcida local pode influenciar, de forma subconsciente, as decisões da
    equipe de arbitragem em lances cruciais, como cartões e penalidades.
\end{itemize}

Essa distinção entre mecanismos é relevante para o presente trabalho. Os quatro
fatores produzem assinaturas empíricas distintas: o viés de arbitragem deve
manifestar-se em cartões e faltas, ao passo que a familiaridade com o ambiente e
a logística afetam diretamente o volume de ações ofensivas. A análise conjunta
dessas métricas permite, portanto, discriminar parcialmente entre os mecanismos,
em vez de tratá-los como um bloco indiviso.

\subsection{O Impacto da Pandemia e os Jogos com Portões Fechados}\label{sec:covid}

Um dos momentos mais singulares para o estudo do Fator Casa ocorreu durante a
pandemia de COVID-19. As restrições sanitárias forçaram a realização de partidas
com portões fechados, conhecidos na Europa como \textit{ghost games}.

Pesquisas como a de \citeonline{Fischer2021} aproveitaram esse cenário de
isolamento para testar empiricamente o peso real da torcida. Isolar essa variável
permitiu comparar o desempenho das equipes em condições anteriormente
impossíveis de observar, o que levou parte da literatura a tratar o episódio como
um experimento natural.

Cabe, entretanto, uma ressalva metodológica que orienta a análise conduzida neste
trabalho. A temporada 2020/21 não se diferenciou das anteriores apenas pela
ausência de público: houve também encurtamento da pré-temporada, congestionamento
de calendário e alteração no número de substituições permitidas por partida.
Comparações entre temporadas, portanto, confundem o efeito da torcida com o
efeito de todas as demais mudanças ocorridas simultaneamente. Contrastes
construídos \textit{dentro} de uma mesma temporada, quando disponíveis, oferecem
evidência mais limpa a respeito do papel específico do público.

\section{A Metodologia CRISP-DM}\label{sec:crisp-dm}

O \textit{Cross Industry Standard Process for Data Mining} (CRISP-DM) é um modelo
de processo que descreve as fases de um projeto de mineração de dados de forma
independente de domínio e de ferramenta~\cite{Wirth2000}. Sua documentação de
referência~\cite{Chapman2000} organiza o trabalho em seis fases:

\begin{enumerate}
  \item \textbf{Entendimento do negócio} (\textit{business understanding}): fase
    inicial focada em compreender os objetivos do projeto sob a perspectiva de
    regras de negócio, transformando esse conhecimento na definição de um
    problema de ciência de dados;
  \item \textbf{Entendimento dos dados} (\textit{data understanding}): envolve a
    coleta inicial dos dados, a identificação de problemas de qualidade e a
    primeira exploração das propriedades do conjunto;
  \item \textbf{Preparação dos dados} (\textit{data preparation}): fase de
    tratamento em que os dados brutos são limpos, transformados, filtrados e
    modelados em conjuntos finais prontos para análise;
  \item \textbf{Modelagem} (\textit{modeling}): seleção e aplicação de técnicas
    estatísticas e algoritmos de aprendizado de máquina, com o ajuste de seus
    parâmetros;
  \item \textbf{Avaliação} (\textit{evaluation}): revisão técnica dos resultados
    obtidos para garantir que as descobertas realmente resolvem o problema
    inicial;
  \item \textbf{Implantação} (\textit{deployment}): organização dos resultados em
    relatórios finais, códigos reaproveitáveis ou \textit{dashboards}.
\end{enumerate}

Uma característica central do CRISP-DM é seu caráter iterativo: as fases não
constituem uma sequência linear, e descobertas feitas em etapas avançadas
frequentemente exigem o retorno a etapas anteriores. Esse aspecto é
particularmente relevante neste trabalho, no qual a fase de modelagem revelou
limitações na preparação dos dados que exigiram revisão retroativa, conforme
descrito no Capítulo~\ref{cap:metodologia}.

\section{Modelos Estatísticos para Resultados de Futebol}\label{sec:modelos-futebol}

Análises descritivas do mando de campo baseiam-se em médias condicionais, como a
média de pontos conquistados em casa e fora. Essa abordagem possui uma limitação
estrutural: ela não controla pela identidade dos adversários enfrentados em cada
condição. O diferencial observado para um clube mistura, de forma
indistinguível, a vantagem efetivamente atribuível ao mando e as características
do calendário que lhe coube.

A literatura estatística do futebol trata esse problema por meio de modelos que
estimam simultaneamente a força de cada equipe e um parâmetro comum de mando.
\citeonline{Maher1982} propôs modelar os gols marcados por cada lado como
variáveis aleatórias de Poisson independentes, com médias determinadas por
parâmetros de ataque e defesa específicos de cada clube.
\citeonline{DixonColes1997} estenderam essa formulação em dois aspectos que se
tornaram referência na área: a introdução de um fator de correção para a
dependência observada nos placares de baixa pontuação, nos quais o modelo de
Poisson independente subestima a frequência de empates, e a ponderação temporal
das observações, de modo que partidas mais recentes influenciem mais a estimativa
da força atual de cada equipe.

Na formulação de Dixon e Coles, para uma partida entre o mandante~$i$ e o
visitante~$j$, os gols são modelados como variáveis de Poisson com médias:

\begin{equation}\label{eq:dixon-coles}
  \lambda_{\text{casa}} = \exp(\mu + \alpha_i - \beta_j + \gamma),
  \qquad
  \lambda_{\text{fora}} = \exp(\mu + \alpha_j - \beta_i),
\end{equation}

\noindent em que $\alpha$ e $\beta$ representam, respectivamente, as forças de
ataque e de defesa de cada clube, $\mu$ é um termo comum de escala e $\gamma$ é o
parâmetro de mando de campo. Como $\gamma$ é estimado conjuntamente com as forças
de todas as equipes, ele mede a vantagem do mandante já descontada a qualidade de
quem enfrentou quem. A quantidade $e^{\gamma} - 1$ expressa o acréscimo
percentual esperado de gols atribuível ao mando.

\section{Avaliação de Modelos Probabilísticos}\label{sec:avaliacao}

A previsão de resultados de futebol é, por natureza, um problema probabilístico:
o interesse não está em afirmar categoricamente qual equipe vencerá, mas em
atribuir probabilidades às três possibilidades de resultado. Essa característica
tem consequências diretas sobre a escolha das métricas de avaliação.

A acurácia, métrica mais intuitiva, é insuficiente nesse contexto por duas
razões. Primeiro, ela ignora a confiança associada a cada previsão, tratando
igualmente um acerto obtido com 51\% de probabilidade e outro obtido com 95\%.
Segundo, em um problema no qual uma das classes é majoritária, a acurácia pode
ser elevada sem que o modelo agregue informação útil.

Por esse motivo, adotam-se regras de pontuação próprias (\textit{strictly proper
scoring rules}), que são minimizadas quando o modelo reporta suas probabilidades
verdadeiras~\cite{GneitingRaftery2007}. As duas mais usuais são o \textit{logarithmic
loss}, definido como a média do logaritmo negativo da probabilidade atribuída ao
resultado observado, e o escore de \citeonline{Brier1950}, definido como o erro
quadrático médio entre as probabilidades previstas e o vetor indicador do
resultado.

Complementarmente, a \textbf{calibração} de um modelo pode ser avaliada
agrupando-se as previsões em faixas de probabilidade e comparando, em cada faixa,
a probabilidade média prevista com a frequência observada do evento. Um modelo
bem calibrado produz pontos próximos à diagonal: entre as partidas às quais
atribui 70\% de chance de vitória do mandante, o mandante deve vencer
aproximadamente 70\% das vezes.

\subsection{Significância Estatística e Tamanho de Efeito}\label{sec:efeito}

Em conjuntos de dados grandes, p-valores reduzidos são facilmente obtidos mesmo
para diferenças de magnitude irrelevante, uma vez que a estatística de teste
cresce com o tamanho da amostra. Por essa razão, recomenda-se que resultados
inferenciais sejam reportados acompanhados de estimativas pontuais, intervalos de
confiança e medidas de tamanho de efeito~\cite{Wasserstein2016}.

Neste trabalho adota-se o $d$ de \citeonline{Cohen1988} como medida de tamanho de
efeito para comparações de médias, interpretado segundo as faixas convencionais
de efeito pequeno, médio e grande, e a razão de chances para comparações de
proporções.

\subsection{Vazamento de Dados e Validação Temporal}\label{sec:vazamento}

Um risco central em problemas preditivos é o vazamento de dados (\textit{data
leakage}), situação em que informação indisponível no momento da previsão é
inadvertidamente incorporada às variáveis explicativas~\cite{Kaufman2012}. Em
previsão de resultados esportivos, o caso mais comum consiste em utilizar
estatísticas da própria partida — como finalizações, escanteios e faltas — que
só passam a existir após o apito final. Modelos construídos dessa forma exibem
desempenho artificialmente elevado e não possuem validade preditiva.

Relacionado a esse risco está o procedimento de validação. Em dados com estrutura
temporal, a partição aleatória entre treino e teste permite que o modelo seja
treinado com observações posteriores às que serão previstas, o que também
superestima o desempenho~\cite{Bergmeir2012}. A alternativa adequada é a
validação temporal, na qual o conjunto de treino é sempre anterior ao conjunto de
teste, reproduzindo as condições reais de aplicação do modelo.
